% Содержимое подраздела 3.3 - Перспективы развития: за пределами RNN и SMILES


\section{Перспективы развития: за пределами RNN и SMILES}
\label{sec:future_perspectives}

Анализ ограничений классических подходов позволяет очертить три магистральных направления, по которым движется современная наука в области генеративного дизайна молекул: переход к более мощным архитектурам, переход к более естественным представлениям данных и, наконец, переход от 2D-структур к 3D-геометрии.

\textbf{Тренд №1: От RNN к Трансформерам и большим языковым моделям.}
Рекуррентные нейронные сети, долгое время доминировавшие в задачах обработки последовательностей, постепенно уступают место архитектуре Трансформер. Благодаря механизму внимания (attention), трансформеры способны улавливать дальние зависимости в последовательностях, что делает их более мощным инструментом для моделирования сложных SMILES-строк. Эта тенденция нашла отражение в современных фреймворках, таких как \textbf{REINVENT 4}, где трансформерная модель была добавлена как опция для задач локальной оптимизации молекул \cite{loeffler2024reinvent}. Последующие работы показали, что обучение с подкреплением (RL) особенно эффективно именно в связке с трансформерами для целенаправленной генерации \cite{he2024evaluation}. Апогеем этого тренда стало появление <<фундаментальных моделей>> (foundation models), таких как \textbf{GP-MOLFORMER}, обученных на миллиардах молекул. Эти модели, заимствуя парадигму из обработки естественного языка (NLP), демонстрируют беспрецедентную универсальность и производительность \cite{ross2025gp}.

\textbf{Тренд №2: От строк к графам.}
Как было показано ранее, строковые представления являются неестественными для молекул, которые по своей природе являются графами. Поэтому вторым ключевым трендом является все более широкое применение графовых нейронных сетей (GNN). Модель \textbf{JT-VAE} была одним из пионеров этого направления. Обширные обзорные работы подтверждают, что GCN (Graph Convolutional Networks) и другие варианты GNN сегодня являются передовой технологией для широкого спектра задач, от предсказания свойств до de novo генерации, так как они способны извлекать признаки напрямую из топологии молекулы, минуя <<бутылочное горлышко>> строкового представления \cite{sun2020graph}.

\textbf{Тренд №3: От 2D к 3D.}
Наиболее передовым и сложным направлением является переход от генерации 2D-графов к генерации полноценных 3D-конформаций молекул. 3D-структура напрямую определяет биологическую активность, и ее моделирование является Святым Граалем вычислительной химии. На этом фронте доминирующим подходом стали \textbf{диффузионные модели}. Такие модели, как \textbf{GeoDiff}, обучаются обращать процесс постепенного зашумления 3D-координат атомов, генерируя физически правдоподобные и геометрически корректные структуры с соблюдением фундаментальных симметрий \cite{xu2022geodiff}. Как и в случае с 2D-моделями, эта бурно развивающаяся область уже столкнулась с необходимостью стандартизации, что привело к появлению комплексных бенчмарков для сравнения 3D-генераторов \cite{qin2025comprehensive}. Финальным аккордом этого тренда можно считать появление моделей вроде \textbf{BindGPT}, которые объединяют в себе все три направления: используют архитектуру трансформера (LLM), работают с единым представлением, объединяющим 2D-граф (SMILES) и 3D-координаты, и интегрируют обучение с подкреплением для целенаправленной генерации молекул с высокой аффинностью к белковой мишени. Это демонстрирует путь к созданию единых, универсальных и мощных инструментов для полностью автоматизированного дизайна лекарств \cite{zholus2025bindgpt}.
